% Rigorous core for Where Fractals Just Touch.
% Standard results and elementary reductions are not claimed as new priority.
% Requires amsmath, amssymb, amsthm; theorem environments supplied by wrapper.
\providecommand{\C}{\mathbb C}
\providecommand{\D}{\mathbb D}
\providecommand{\cK}{\mathcal K}
\providecommand{\cM}{\mathcal M}
\providecommand{\cW}{\mathcal W}
\providecommand{\dist}{\operatorname{dist}}
\providecommand{\diam}{\operatorname{diam}}
\providecommand{\Fix}{\operatorname{Fix}}
\providecommand{\Ker}{\operatorname{ker}}

\section{Canonical normalization and the atlas coordinates}
\label{sec:core-normalization}

Put $J_+(z)=z$ and $J_-(z)=\overline z$. A contracting planar
similarity is $F(z)=c+\lambda J_\kappa(z)$, where
$\kappa\in\{+,-\}$ and $0<|\lambda|<1$. Its plane orientation is
$\kappa$. This discrete datum must be distinguished from the endpoint
reversal signature of a zipper, introduced below.

\begin{proposition}[Canonical translation normalization]
\label{prop:canonical-normalization}
Let $F_1,F_2$ be contracting planar similarities with distinct fixed
points. There is a unique complex-affine coordinate change
$H(z)=Az+B$, with $A\ne0$, for which
\begin{equation}\label{eq:canonical-system}
H F_1 H^{-1}(z)=-1+aJ_{\kappa_1}(z),\qquad
H F_2 H^{-1}(z)=1+bJ_{\kappa_2}(z).
\end{equation}
The coefficients satisfy $0<|a|,|b|<1$. Consequently each labelled
orientation class has the parameter domain
$\mathcal P=(\D\setminus\{0\})^2$, of real dimension four.
Pairs with a common fixed point have a singleton attractor and are
excluded from this normalization.
\end{proposition}

\begin{proof}
The real-affine map $M=(F_1+F_2)/2$ is a strict contraction, with
Lipschitz constant at most $(r_1+r_2)/2<1$. Let $c$ be its unique
fixed point and put
\[
d=F_2(c)-c=-(F_1(c)-c).
\]
If $d=0$, both similarities fix $c$, contrary to the hypothesis.
Thus $H(z)=(z-c)/d$ is an admissible coordinate change and its
conjugated maps have the stated constant terms. A direct coefficient
$\lambda$ is unchanged; an opposite coefficient becomes
$\lambda\overline d/d$, so contraction ratios and plane orientations
are preserved.

Conversely, any such normalization has to send the fixed point of
$M$ to the fixed point $0$ of the normalized average. It must also
send $F_2(c)$ to $1$. These two conditions determine $H$ uniquely.
If the original maps have a common fixed point, that singleton is
their compact invariant set, hence their attractor by uniqueness
\cite{Hutchinson1981}. Finally, the normalized maps cannot have a
common fixed point: such a point would be fixed by their average,
hence would be $0$, whereas their images of $0$ are $-1$ and $1$.
\end{proof}

We also write $S_-=S_1$ and $S_+=S_2$ for the two normalized maps.
Relabelling identifies the mixed classes $(+,-)$ and $(-,+)$,
leaving three unordered plane-orientation types $(+,+),(+,-),(-,-)$.
This is a classification of the discrete types, not an assertion
that there is only one affine-conjugacy class within each type.
In the normalized chart, relabelling followed by $z\mapsto-z$
exchanges $(a,b)$ and the orientation labels; complex conjugation
replaces $(a,b)$ by $(\overline a,\overline b)$ and preserves the
orientation labels. The normalization removes any further continuous
freedom from complex-affine conjugacy. General real-affine
coordinate changes, which may include shears, are not part of
this equivalence relation.

Write $r_1=|a|$, $r_2=|b|$. The coordinates used by the atlas are
\begin{equation}\label{eq:atlas-core-coordinates}
w=\frac{r_1}{r_1+r_2},\qquad \rho=\frac{2}{r_1+r_2},\qquad
a=\frac{2w}{\rho}e^{i(\alpha-\gamma)},\quad
b=\frac{2(1-w)}{\rho}e^{-i(\alpha+\gamma)}.
\end{equation}
Here $0<w<1$ and the strict contraction condition is
$\rho>2\max\{w,1-w\}$. The phase variables belong to a torus:
\begin{equation}\label{eq:phase-lattice}
(\alpha,\gamma)\in\mathbb R^2/\Lambda,\qquad
\Lambda=\langle(\pi,\pi),(\pi,-\pi)\rangle_{\mathbb Z}.
\end{equation}
Indeed, if $\theta_1,\theta_2$ are arguments of $a,b$, then
$\alpha=(\theta_1-\theta_2)/2$ and
$\gamma=-(\theta_1+\theta_2)/2$. One can use
$0\le\alpha\le\pi$, with $\gamma$ modulo $2\pi$, provided the seam
$(0,\gamma)\sim(\pi,\gamma+\pi)$ is retained. Endpoints of this
rectangle are not independent parameter values.

At fixed $w$, the remaining parameter domain is three-dimensional. A display
of sampled radial transitions therefore provides a way to navigate this
section of the atlas. The angular identifications in~\eqref{eq:phase-lattice}
and the sampled radial intervals are retained with each point. Section~\ref{sec:explorer-methods}
specifies the display map and its relation to the underlying parameters.


